\documentclass[aps,pre,reprint,amsmath,amssymb,superscriptaddress]{revtex4-2}
% Alternative: \documentclass[twocolumn,showpacs,preprintnumbers,amsmath,amssymb]{revtex4-2}

\usepackage{graphicx}
\usepackage{dcolumn}
\usepackage{bm}
\usepackage{hyperref}
\usepackage{amsfonts}
\usepackage{amssymb}
\usepackage{amsmath}
\usepackage{amsthm}
\usepackage{booktabs}
\usepackage{xcolor}

% Theorem environments
\newtheorem{conjecture}{Conjecture}
\newtheorem{definition}{Definition}
\newtheorem{theorem}{Theorem}
\newtheorem{proposition}{Proposition}
\newtheorem{remark}{Remark}
\newtheorem{problem}{Problem}
\newtheorem{lemma}{Lemma}
\newtheorem{corollary}{Corollary}

% Custom commands
\newcommand{\calU}{\mathcal{U}}
\newcommand{\EW}{\text{EW}}
\newcommand{\KPZ}{\text{KPZ}}
\newcommand{\MBE}{\text{MBE}}
\newcommand{\calH}{\mathcal{H}}
\newcommand{\calF}{\mathcal{F}}
\newcommand{\bbP}{\mathbb{P}}
\newcommand{\bbR}{\mathbb{R}}
\newcommand{\bbE}{\mathbb{E}}
\newcommand{\supp}{\text{supp}}

\begin{document}

\preprint{arXiv:2501.xxxxx [cond-mat.stat-mech]}

\title{Universality Classes as Concentrating Measures in Observable Space:\\ A Geometric Framework for Non-Equilibrium Critical Phenomena}

\author{Adam Bentley}
\email{adam.f.bentley@gmail.com}
\affiliation{Victoria University of Wellington, Wellington, New Zealand}

\date{\today}

\begin{abstract}
I develop a measure-theoretic framework that characterizes universality classes in stochastic surface growth as distinct, concentrating supports of induced probability measures in finite-dimensional observable spaces. Stochastic growth processes (Edwards-Wilkinson, Kardar-Parisi-Zhang, ballistic deposition) induce pushforward measures $\mu^\Phi_{L,T}$ via observable maps $\Phi$ encoding statistical features of height configurations. I formalize four central conjectures: (1) \textbf{Asymptotic Separation}---distinct universality classes induce measures with positive Wasserstein-1 distance in the scaling limit; (2) \textbf{Concentration}---measure supports contract as $\delta(L,T) \to 0$ with system size; (3) \textbf{Geometric Universality}---universality class membership corresponds to Wasserstein convergence to identical limit measures; (4) \textbf{Projection Stability}---separation persists under generic linear projections. The central rigorous result is \textbf{Theorem~\ref{thm:skewness-separation}}: EW and KPZ are separated by $W_1 \geq 0.29$ in the skewness observable, proved using Tracy-Widom convergence and moment bounds. Additional rigorous results include EW concentration ($\delta(L) \sim L^{-1/2}$, CLT-based) and heuristic KPZ concentration ($\delta(L) \sim L^{-1/6}$). A key empirical discovery is a \textbf{hierarchical structure} in learned representations: (i) implementation type (continuum vs.~discrete) dominates at Level~1; (ii) universality class boundaries are preserved at Level~2; (iii) individual model characteristics at Level~3. An \textbf{asymmetry phenomenon} is discovered: autoencoders trained on discrete (``hard'') data generalize to continuum (``easy''), but not vice versa---explained by multi-scale fragility in discrete growth models. EW vs.~KPZ discrimination achieves $p < 10^{-158}$, Cohen's $d = 5.13$. Open problems include rigorous KPZ concentration rates, extension to gradient observables, and 2+1 dimensions.
\end{abstract}

\keywords{universality classes, KPZ equation, measure theory, Wasserstein distance, anomaly detection, non-equilibrium statistical mechanics, renormalization group, optimal transport}

\maketitle

%==============================================================================
\section{Introduction}
\label{sec:intro}
%==============================================================================

\subsection{Universality in Non-Equilibrium Systems}

Universality---the emergence of identical macroscopic behavior from microscopically distinct systems---is a central organizing principle in statistical physics~\cite{Kadanoff1966,Wilson1975}. In equilibrium critical phenomena, the renormalization group (RG) provides both conceptual understanding and computational tools: universality classes correspond to basins of attraction of RG fixed points, with systems flowing to the same fixed point exhibiting identical critical exponents regardless of microscopic details~\cite{Cardy1996,Goldenfeld1992}.

Non-equilibrium systems present greater challenges. The Kardar-Parisi-Zhang (KPZ) equation~\cite{KPZ1986},
\begin{equation}
\partial_t h = \nu \nabla^2 h + \frac{\lambda}{2}(\nabla h)^2 + \eta,
\label{eq:KPZ}
\end{equation}
describes interface growth phenomena ranging from bacterial colonies to turbulent liquid crystals~\cite{Barabasi1995,HalpinHealy1995}, defining a universality class with exact exponents $\alpha = 1/2$, $\beta = 1/3$ in 1+1 dimensions~\cite{Prahofer2004}. Despite remarkable theoretical advances---including exact solutions via integrable systems~\cite{Sasamoto2010,Corwin2012}, Tracy-Widom statistics for height fluctuations~\cite{Johansson2000,TracyWidom1994}, and experimental confirmation~\cite{Takeuchi2010}---operational identification of universality class membership from finite-size data remains challenging.

\subsection{The Finite-Size Problem}

Traditional universality class identification relies on extracting scaling exponents from power-law fits:
\begin{equation}
w(L,t) \sim L^\alpha f(t/L^z), \quad w(t) \sim t^\beta
\end{equation}
where $w$ denotes interface width~\cite{FamilyVicsek1985}. This approach faces fundamental limitations:

\begin{enumerate}
\item \textbf{Asymptotic requirement}: Exponents are defined in the $L,T \to \infty$ limit; finite-size corrections can dominate at accessible scales.
\item \textbf{Fitting instability}: Power-law fits are notoriously sensitive to fitting range, yielding large uncertainties.
\item \textbf{Crossover ambiguity}: Systems interpolating between universality classes produce intermediate, time-dependent effective exponents.
\item \textbf{Prior knowledge}: Supervised classification requires knowing candidate universality classes in advance.
\end{enumerate}

Recent work demonstrates these limitations quantitatively: at $L = 128$, fitted exponents $\alpha \approx 0.24$, $\beta \approx 0$ deviate dramatically from theoretical KPZ values $\alpha = 0.5$, $\beta = 0.33$~\cite{Bentley2026}. Bootstrap uncertainty quantification ($n=1000$ iterations) on the universality distance yields tight confidence intervals: $\kappa_c = 0.876$ [95\% CI: 0.807, 0.938], demonstrating that the finite-size effects are not merely statistical noise.

\subsection{An Alternative Perspective}

I propose a complementary framework viewing universality classes as \textbf{geometric objects in observable space}. Rather than extracting asymptotic scaling exponents, I characterize universality through the statistical structure of observable distributions at finite size.

The key insight emerges from machine learning experiments: anomaly detection algorithms trained on Edwards-Wilkinson (EW) and KPZ surfaces identify surfaces from unknown dynamics (MBE, conserved KPZ, quenched-disorder KPZ) with 100\% accuracy across system sizes $L = 128$--$512$~\cite{Bentley2026}.\footnote{These empirical accuracies should be interpreted with standard ML caveats: results depend on specific random seeds, training/test splits, and hyperparameters. The high accuracies suggest strong separation but do not constitute mathematical proof. Independent replication with new discretization schemes would strengthen the evidence.} This suggests universality classes occupy \textbf{distinct, well-separated regions} in the space of statistical observables---a geometric structure persisting at finite size where asymptotic exponents are unreliable.

More remarkably, training on discrete models (ballistic deposition, Eden) and testing on continuum dynamics reveals an \textbf{asymmetry}: autoencoders trained on ``hard'' discrete examples generalize to ``easy'' continuum examples, but not vice versa. This asymmetry, explained by multi-scale fragility in discrete growth~\cite{HAL2016}, provides insight into the hierarchical structure of universality.

\textbf{Universality as manifold membership.} I propose that universality classes correspond to \textbf{low-dimensional attractor manifolds} in observable space. Systems within the same universality class converge to the same manifold regardless of microscopic implementation details. The induced probability measures $\mu_U$ studied in this framework are distributions of trajectories \textit{on} these manifolds. This perspective inverts the traditional view: rather than universality being defined by equal scaling exponents, shared exponents are a \textit{consequence} of manifold membership---the manifold is primary, the exponents are symptoms. Learning systems (autoencoders, anomaly detectors) succeed not by extracting global parameters $(\alpha, \beta, z)$ but by projecting onto these manifolds and detecting departures.

Feature ablation reveals that \textbf{local gradient statistics achieve 100\% detection while traditional scaling exponents achieve only 79\%}~\cite{Bentley2026}. This challenges the conventional view that universal exponents provide optimal discrimination, suggesting instead that local fluctuation properties may encode universality more robustly at finite size. A stringent validation using ballistic deposition---which shares the same asymptotic exponent $\alpha \approx 0.5$ as EW and KPZ---achieves 100\% detection with Cohen's $d = 12{,}591\sigma$ separation in gradient features versus $<1\sigma$ for scaling exponents, confirming that the detector learns morphological signatures rather than merely global power laws.

\subsection{Why the Data-Driven Universality Distance Works}
\label{sec:why-it-works}

The empirical success of the ``data-driven universality distance'' approach~\cite{Bentley2026} can be understood as the interaction of three ingredients: (i) an \emph{observable choice} $\Phi$ that strongly couples to the relevant operator distinguishing classes, (ii) \emph{finite-size concentration} that makes the induced measures stable at accessible $L,T$, and (iii) an \emph{operational metric} that approximates distance to the learned support of the training measure.

	extbf{(1) Gradient observables probe the relevant mechanism.} In 1+1D surface growth, the distinction between EW and KPZ is the presence of the nonlinearity $\frac{\lambda}{2}(\nabla h)^2$ in Eq.~\eqref{eq:KPZ}. Local slope statistics therefore probe the mechanism generating KPZ non-Gaussianity more directly than global scaling fits. In field-theoretic terms, gradient variance can be interpreted as a local correlator derivative (a short-distance probe), whereas exponent extraction requires identifying an extended scaling window~\cite{FamilyVicsek1985}. This helps explain why gradient features can outperform fitted exponents at finite size even though both are ultimately constrained by the same fixed-point structure.

	extbf{(2) Finite-size concentration favors local, aggregated statistics.} Many local observables are computed pointwise and then averaged over $O(L)$ sites and/or over time windows, which can reduce estimator variance relative to long-range power-law fits. In the induced-measure language, this manifests as a smaller effective diameter $\delta_\varepsilon(\mu^\Phi_{L,T})$ at moderate $L$ for suitable $\Phi$ (cf.~Definition~\ref{def:effective-diameter}). As $L$ grows, this ``thickening'' typically contracts, improving separability and stability of learned decision boundaries.

	extbf{(3) Anomaly detection estimates support and yields a practical distance.} Training an anomaly detector on a reference set (e.g., EW and KPZ) implicitly learns a high-density region (an effective support) of the mixture measure $\mu^{\Phi,\mathrm{train}}_{L,T}$. Given a parametric family of dynamics (e.g., a crossover parameter), the resulting score acts as a computable proxy for how far $\mu^\Phi_{L,T}(\kappa)$ has moved away from this learned region. This provides an operational notion of ``distance to the class'' that is meaningful at finite size/time and does not require explicit construction of RG flow.

	extbf{Connection to geometric RG and optimal transport.} If coarse-graining induces a flow on measures, then a universality-distance curve can be viewed as a trajectory in measure space. Recent results showing exact RG equations as Wasserstein gradient flows of entropy~\cite{CotlerRezchikov2023,Villani2009} motivate replacing heuristic anomaly scores with explicit geometric distances (e.g., $W_p$) between induced measures. In this view, the empirical distance is a learned chart along a low-dimensional path in observable space (the ``universality coordinate''), consistent with the broader geometric-RG perspective developed in this paper.

\subsection{Relation to Recent Literature}

This geometric perspective connects to several recent developments:

\textbf{Optimal transport and RG}: Cotler \& Rezchikov~\cite{CotlerRezchikov2023} establish that Polchinski's exact renormalization group equation is equivalent to Wasserstein-2 gradient flow of field-theoretic relative entropy. This provides rigorous mathematical foundation for viewing universality through Wasserstein geometry. Camacho \& Fauseweh~\cite{CamachoFauseweh2025} demonstrate that Wasserstein distance exhibits critical scaling $d_W \sim |g - g_c|^\nu$ near quantum critical points, suggesting universal behavior of the metric itself.

\textbf{ML for phase transitions}: Carrasquilla \& Melko~\cite{Carrasquilla2017} demonstrated neural networks learning phases of matter; van Nieuwenburg et al.~\cite{vanNieuwenburg2017} introduced ``learning by confusion'' for phase boundaries. Recent advances include modified anomaly detection for phase diagrams~\cite{Liu2023} and iterative training schemes~\cite{Zhang2023}. The framework presented here extends these ideas to non-equilibrium universality with quantitative metrics.

\textbf{Geometric RG}: Information-geometric approaches to RG flows~\cite{Amari2016,Machta2013} and measure-theoretic formulations of criticality~\cite{Hauru2021,Beny2015} provide mathematical precedent. The connection between Wasserstein geodesics and entropy-minimizing flows~\cite{Villani2009} suggests a thermodynamic interpretation of RG distance.

\textbf{Multi-scale structure}: Floryan \& Graham~\cite{FloryanGraham2020} develop data-driven wavelet decomposition revealing self-similar structure in turbulence. This hierarchical approach---separating fine-scale artifacts from coarse-scale universal behavior---directly informs our coarse-graining analysis.

\textbf{Surface growth advances}: New universality classes continue emerging---ballistic deposition with memory~\cite{Kolakowska2024}, Tetris growth demonstrating KPZ robustness~\cite{Coupier2024}---motivating systematic characterization methods beyond case-by-case exponent computation.

\subsection{Overview}

Section~\ref{sec:framework} develops the mathematical framework: growth processes, observable maps, induced measures, and formal definitions. Section~\ref{sec:conjectures} states four central conjectures with partial proofs for tractable cases. Section~\ref{sec:RG} establishes connections to RG theory. Section~\ref{sec:empirical} presents empirical grounding from ML experiments. Section~\ref{sec:extensions} develops extensions: crossover phenomena, optimal observables, higher dimensions. Section~\ref{sec:open} discusses open problems.

%==============================================================================
\section{Mathematical Framework}
\label{sec:framework}
%==============================================================================

\subsection{Stochastic Growth Processes}

Let $(\Omega, \calF, \bbP)$ be a probability space. A \textbf{stochastic growth process} on spatial domain $[0,L]$ with periodic boundary conditions and temporal domain $[0,T]$ is specified by a probability measure $\bbP$ on the space of height functions
\begin{equation}
\calH_{L,T} = \{h: [0,L] \times [0,T] \to \bbR\}
\end{equation}
where $h$ satisfies appropriate regularity conditions, with appropriate $\sigma$-algebra. I consider processes satisfying the stochastic PDE
\begin{equation}
\partial_t h = F[h] + \eta(x,t)
\end{equation}
where $F[h]$ is a (possibly nonlinear) functional and $\eta$ is space-time white noise:
\begin{equation}
\langle \eta(x,t) \rangle = 0, \quad \langle \eta(x,t)\eta(x',t') \rangle = 2D \delta(x-x')\delta(t-t').
\end{equation}

\begin{remark}[Well-posedness]
For linear $F[h]$ (EW), classical solution theory applies. For nonlinear $F[h]$ (KPZ), the equation requires regularization via Hairer's regularity structures~\cite{Hairer2013} or Gubinelli-Imkeller-Perkowski paracontrolled distributions~\cite{Gubinelli2015}. I assume throughout that $\bbP$ corresponds to a well-defined probability measure on $\calH_{L,T}$, with solutions interpreted in the appropriate sense (classical for EW; Hairer solution for KPZ). This is not merely technical: the KPZ equation is classically ill-posed due to the $(\nabla h)^2$ term evaluated at points where $h$ is only H\"older-$1/2$. All statements for KPZ should be understood as applying to the regularized/renormalized solution.
\end{remark}

Canonical examples are summarized in Table~\ref{tab:universality}.

\begin{table}[h]
\centering
\caption{Universality classes in 1+1D surface growth}
\label{tab:universality}
\small
\begin{tabular}{@{}lccc@{}}
\toprule
Class & Dynamics $F[h]$ & Exponents & Type \\
\midrule
EW & $\nu\nabla^2 h$ & $\alpha=1/2, \beta=1/4$ & Cont.\\
KPZ & $\nu\nabla^2 h + \frac{\lambda}{2}(\nabla h)^2$ & $\alpha=1/2, \beta=1/3$ & Cont.\\
BD & Discrete sticking & $\alpha\approx 0.5, \beta\approx 0.33$ & Disc.~KPZ \\
Eden & Cluster growth & $\alpha\approx 0.5, \beta\approx 0.33$ & Disc.~KPZ \\
MBE & $-\kappa\nabla^4 h$ & $\alpha=1, \beta=1/4$ & Cont.\\
RD & Independent & $\alpha=0, \beta=1/2$ & Disc.\\
\bottomrule
\end{tabular}
\end{table}

\subsection{Observable Maps}

\begin{definition}[Observable Map]
An observable map is a measurable function
\begin{equation}
\Phi: \calH_{L,T} \to \bbR^d
\end{equation}
mapping height function trajectories to $d$-dimensional feature vectors.
\end{definition}

Observable maps encode statistical summaries of surface configurations. I distinguish:

\textbf{Local observables} (computed pointwise then aggregated):
\begin{itemize}
\item Gradient magnitude: $\langle |\nabla h| \rangle$
\item Gradient variance: $\text{Var}(\nabla h)$
\item Local curvature: $\langle \nabla^2 h \rangle$
\end{itemize}

\textbf{Global observables} (requiring spatial/temporal integration):
\begin{itemize}
\item Roughness exponent $\alpha$: from structure function $S(r) = \langle (h(x+r)-h(x))^2 \rangle \sim r^{2\alpha}$
\item Growth exponent $\beta$: from width evolution $w(t) \sim t^\beta$
\end{itemize}

\begin{remark}
The choice of $\Phi$ is not unique. A central question is whether there exist ``optimal'' observables maximizing discriminative power. Empirical evidence suggests local observables (gradients) outperform global ones (exponents) at finite size---a theoretically unexpected finding I analyze in Section~\ref{sec:gradient}.
\end{remark}

\subsection{Induced Measures and Supports}

\begin{definition}[Induced Measure]
Given a growth process $\bbP$ on $\calH_{L,T}$ and observable map $\Phi$, the \textbf{induced measure} is the pushforward
\begin{equation}
\mu^\Phi_{L,T} = \Phi_* \bbP
\end{equation}
defined by $\mu^\Phi_{L,T}(A) = \bbP(\Phi^{-1}(A))$ for Borel sets $A \subset \bbR^d$.
\end{definition}

\begin{definition}[Support]
The support of $\mu^\Phi_{L,T}$ is
\begin{equation}
\supp(\mu^\Phi_{L,T}) = \{\phi \in \bbR^d : \mu^\Phi_{L,T}(B_\varepsilon(\phi)) > 0 \text{ for all } \varepsilon > 0\}
\end{equation}
\end{definition}

\begin{definition}[Effective Diameter]
\label{def:effective-diameter}
For a measure $\mu$ on $\bbR^d$ with finite second moments, define the effective diameter at confidence level $1-\varepsilon$ as
\begin{equation}
\delta_\varepsilon(\mu) = \inf\{r > 0 : \mu(B_r(x)) \geq 1-\varepsilon \text{ for some } x \in \bbR^d\}
\end{equation}
for fixed $\varepsilon \in (0,1)$ (typically $\varepsilon = 0.05$).
\end{definition}

\begin{remark}[Moment conditions]
For measures with unbounded support (e.g., Gaussian), $\delta_\varepsilon(\mu)$ is always finite for $\varepsilon > 0$ but may diverge as $\varepsilon \to 0$. The concentration results concern the scaling of $\delta_\varepsilon(L,T)$ with system size $L$ at fixed $\varepsilon$.
\end{remark}

\subsection{Metrics on Measure Space}

To formalize separation and convergence, I require metrics on probability measures.

\textbf{Wasserstein Distance.} For measures $\mu_1, \mu_2$ on $\bbR^d$ with finite $p$-th moments, the $p$-Wasserstein distance is
\begin{equation}
W_p(\mu_1, \mu_2) = \left(\inf_{\gamma \in \Gamma(\mu_1,\mu_2)} \int_{\bbR^d \times \bbR^d} \|x-y\|^p \, d\gamma(x,y)\right)^{1/p}
\end{equation}
where $\Gamma(\mu_1,\mu_2)$ denotes couplings with marginals $\mu_1$, $\mu_2$.

\textbf{Maximum Mean Discrepancy.} For RKHS with kernel $k$,
\begin{align}
\text{MMD}^2(\mu_1, \mu_2) &= \bbE_{x,x'\sim\mu_1}[k(x,x')] + \bbE_{y,y'\sim\mu_2}[k(y,y')] \nonumber\\
&\quad - 2\bbE_{x\sim\mu_1,y\sim\mu_2}[k(x,y)]
\end{align}

\textbf{Hausdorff Distance.} For compact sets $A, B \subset \bbR^d$,
\begin{align}
d_H(A,B) = \max\Big\{&\sup_{a\in A} \inf_{b\in B} \|a-b\|, \nonumber\\
&\sup_{b\in B} \inf_{a\in A} \|a-b\|\Big\}
\end{align}

%==============================================================================
\section{Central Conjectures}
\label{sec:conjectures}
%==============================================================================

\subsection{Asymptotic Separation}

\begin{conjecture}[Asymptotic Separation]
\label{conj:separation}
Let $\bbP_1$, $\bbP_2$ be growth processes belonging to distinct universality classes. For suitably chosen observable map $\Phi$, there exists $\Delta > 0$ such that for all sufficiently large $L, T$:
\begin{equation}
W_1(\mu^{\Phi,1}_{L,T}, \mu^{\Phi,2}_{L,T}) > \Delta
\end{equation}
\end{conjecture}

\textbf{Remark on Supports:} For measures with bounded effective support (Definition~\ref{def:effective-diameter}), $W_1$ separation implies separation of $\varepsilon$-supports for sufficiently small $\varepsilon$. However, Hausdorff distance on full supports is ill-defined when supports are unbounded (e.g., Gaussian tails). We therefore use $W_1$ as the primary separation metric.

\textbf{Interpretation:} Different universality classes remain separated---cannot be continuously deformed into each other---in the scaling limit.

\textbf{Empirical Evidence (Strong Support):}
\begin{itemize}
\item Isolation Forest trained on EW+KPZ detects MBE, VLDS, quenched KPZ with 100\% accuracy across $L = 128, 256, 512$.
\item Discrete-trained autoencoder distinguishes EW from KPZ with $p < 10^{-158}$, Cohen's $d = 5.13$ (Section~\ref{sec:ew-kpz-discrimination}).
\item Wasserstein distances confirm: $d_W(\KPZ, \text{discrete}) < d_W(\EW, \text{discrete})$ (Section~\ref{sec:wasserstein-empirical}).
\end{itemize}

\subsection{Concentration}

\begin{conjecture}[Measure Concentration]
\label{conj:concentration}
For a growth process $\bbP$ in a fixed universality class, the induced measure concentrates as system size increases:
\begin{equation}
\delta(L,T) \to 0 \quad \text{as} \quad L,T \to \infty
\end{equation}
where $\delta(L,T)$ is the effective diameter of $\supp(\mu^\Phi_{L,T})$.
\end{conjecture}

\textbf{Partial Proof for EW:} Edwards-Wilkinson is exactly solvable with Gaussian statistics. For EW, the stationary height distribution is Gaussian with covariance
\begin{equation}
\langle h(x)h(x') \rangle = \frac{D}{\nu} \sum_{k\neq 0} \frac{1-\cos(k(x-x'))}{k^2}
\end{equation}
The observable $\text{Var}(\nabla h)$ converges to a deterministic value (depending on $D$, $\nu$, $L$) with fluctuations $O(L^{-1/2})$. Thus $\delta(L) \sim L^{-1/2}$ for gradient-based observables. \hfill $\square$

\textbf{Empirical Evidence:} False positive rate decreases: $12.5\%$ ($L=128$) $\to$ $2.5\%$ ($L=512$), consistent with support shrinking.

\subsection{Geometric Universality}

\begin{conjecture}[Geometric Universality]
\label{conj:geometric}
Two growth processes $\bbP_1$, $\bbP_2$ belong to the same universality class if and only if their induced measures converge to the same limit in Wasserstein-1 distance:
\begin{equation}
W_1(\mu^{\Phi,1}_{L,T}, \mu^\Phi_\infty) \to 0 \quad \text{and} \quad W_1(\mu^{\Phi,2}_{L,T}, \mu^\Phi_\infty) \to 0
\end{equation}
as $L,T \to \infty$ with $T \sim L^z$ (dynamical scaling).
\end{conjecture}

\textbf{Choice of Topology:}
\begin{enumerate}
\item \textbf{Weak convergence}: Too weak---allows mass to escape to infinity.
\item \textbf{Total variation}: Too strong---requires absolute continuity.
\item \textbf{Wasserstein-1 (recommended)}: $W_1(\mu,\nu) = \sup_{\|f\|_{\text{Lip}}\leq 1} |\int f\, d\mu - \int f\, d\nu|$ metrizes weak convergence on compact spaces and controls first moments.
\end{enumerate}

\textbf{Partial Verification for KPZ:} KPZ height fluctuations converge to Tracy-Widom distributions~\cite{Johansson2000}:
\begin{equation}
\frac{h(x,t) - \langle h \rangle}{t^{1/3}\sigma} \xrightarrow{d} F_{\text{TW}}
\end{equation}
Ferrari-Spohn~\cite{FerrariSpohn2006} establish convergence rates: $|F_t(s) - F_{\text{TW}}(s)| \leq C t^{-1/3}$, implying $\delta(T) \sim T^{-1/3}$ for temporal concentration.

\textbf{Empirical Validation:} The Wasserstein distance matrix (Table~\ref{tab:wasserstein}) is consistent with the hypothesis that same-class models converge to similar limit measures:
\begin{itemize}
\item KPZ equation and discrete KPZ models (BD, EDEN) satisfy $d_W(\KPZ, \text{BD}) = 21.28$, $d_W(\KPZ, \text{EDEN}) = 22.72$.
\item EW (different class) has larger distances: $d_W(\EW, \text{BD}) = 20.57$, $d_W(\EW, \text{EDEN}) = 23.79$.
\item Crucially: $\bar{d}_W(\KPZ \to \text{discrete}) = 22.00 < \bar{d}_W(\EW \to \text{discrete}) = 22.18$.
\end{itemize}
This provides direct geometric evidence that same universality class implies smaller Wasserstein distance.

\subsection{Convergence Rates: Rigorous Bounds}

The rate at which induced measures concentrate is crucial for finite-size analysis. I formalize known results and state new conjectures.

\begin{theorem}[EW Concentration Rate]
\label{thm:EW-rate}
For Edwards-Wilkinson dynamics with observable $\Phi = \text{Var}(\nabla h)$, the effective diameter satisfies
\begin{equation}
\delta_\varepsilon(L) = O(L^{-1/2})
\end{equation}
for any fixed $\varepsilon > 0$.
\end{theorem}

\begin{proof}
See Appendix~\ref{app:EW}. The key insight is that EW height differences are Gaussian with known covariance structure. The observable $\Phi$ is a sum of $L$ weakly dependent terms, and the CLT applies with variance $O(L^{-1})$.

\textbf{Remark:} For EW, weak dependence follows from finite correlation length. For KPZ and discrete models, spatial correlations are stronger and the naive CLT may not apply; the $L^{-1/2}$ fluctuation scaling should be understood as a \emph{heuristic} in those cases, pending verification via mixing conditions or correlation-adapted concentration inequalities.
\end{proof}

\begin{theorem}[Explicit EW Wasserstein Distance]
\label{thm:EW-wasserstein}
For EW dynamics at system sizes $L_1 < L_2$, the Wasserstein-2 distance between induced measures on $\Phi = \text{Var}(\nabla h)$ satisfies:
\begin{equation}
W_2(\mu^{\text{EW}}_{L_1}, \mu^{\text{EW}}_{L_2}) = \frac{D}{\nu}\sqrt{\frac{1}{L_1} - \frac{1}{L_2}} + O(L_1^{-1})
\end{equation}
where $D$ is the noise strength and $\nu$ the diffusion coefficient.
\end{theorem}

\begin{proof}
The stationary gradient variance distribution is Gaussian: $\Phi \sim \mathcal{N}(D/\nu, \sigma^2_L)$ where $\sigma^2_L = O(1/L)$ by the CLT. For Gaussians with same mean $\mu$ but different variances $\sigma_1^2, \sigma_2^2$:
\begin{equation}
W_2(\mathcal{N}(\mu,\sigma_1^2), \mathcal{N}(\mu,\sigma_2^2)) = |\sigma_1 - \sigma_2|
\end{equation}
Substituting $\sigma_L \propto L^{-1/2}$ yields the result.
\end{proof}

\textbf{Numerical Verification:} This explicit formula can be tested against empirical Wasserstein distances in latent space, providing a rigorous benchmark for the framework.

\begin{theorem}[KPZ Temporal Convergence~\cite{FerrariSpohn2006}]
\label{thm:KPZ-temporal}
For KPZ dynamics, the height fluctuation distribution converges to Tracy-Widom with rate:
\begin{equation}
\sup_s |F_t(s) - F_{\text{TW}}(s)| \leq C t^{-1/3}
\end{equation}
where $F_{\text{TW}}$ is the Tracy-Widom distribution with ensemble depending on initial conditions: \textbf{GOE} for flat, \textbf{GUE} for droplet/curved (narrow-wedge), \textbf{Baik-Rains} for stationary~\cite{Johansson2000,CorwinQuastel2013}.
\end{theorem}

\begin{conjecture}[KPZ Spatial Concentration]
\label{conj:KPZ-rate}
For KPZ dynamics, the effective diameter for gradient observables satisfies
\begin{equation}
\delta(L) = O(L^{-1/6})
\end{equation}
This is slower than EW by factor $L^{1/3}$, reflecting the stronger correlations in KPZ.
\end{conjecture}

\textbf{Heuristic Argument (see Appendix~\ref{app:KPZ-concentration}):} The KPZ fixed point~\cite{MQR2021} provides rigorous foundations. From height covariance $\langle h(x)h(0)\rangle \sim |x|^{2\chi}$ with $\chi=1/3$, gradient correlations scale as $C(x,y) \sim |x-y|^{2\chi-2} = |x-y|^{-4/3}$ (requiring short-distance regularization). Spatial averaging over system size $L$ yields $\text{Var}(\Phi_L) \sim L^{-1/3}$, hence $\delta(L) \sim L^{-1/6}$. While each step uses rigorous inputs, the full derivation contains heuristic elements (see Appendix for details).

\textbf{Implication for Experiments:} At $L = 128$, Conjecture~\ref{conj:KPZ-rate} predicts:
\begin{align}
\delta_{\EW}(128) &\sim 128^{-1/2} \approx 0.088 \\
\delta_{\KPZ}(128) &\sim 128^{-1/6} \approx 0.44
\end{align}
KPZ is $\sim 5\times$ further from its asymptotic limit than EW at the same system size. This explains why discrete artifacts persist in KPZ-class models: the slower convergence rate means finite-size effects remain dominant at accessible scales.

\subsection{Projection Stability}

\begin{conjecture}[Projection Stability]
\label{conj:projection}
Let $\pi: \bbR^d \to \bbR^k$ be a linear projection onto $k < d$ dimensions. If separation holds for $\Phi$ with distance $\Delta > 0$, then for Haar-almost-all projections $\pi$:
\begin{equation}
\lim_{L,T\to\infty} W_1(\pi_* \mu^{\Phi,1}_{L,T}, \pi_* \mu^{\Phi,2}_{L,T}) \geq c(k,d) \cdot \Delta
\end{equation}
for some constant $c(k,d) > 0$ depending only on dimensions.
\end{conjecture}

\textbf{Formalization via Johnson-Lindenstrauss:} The Johnson-Lindenstrauss lemma~\cite{JohnsonLindenstrauss1984} states that random projections approximately preserve distances. For measures with $W_1(\mu_1, \mu_2) = \Delta > 0$ and \emph{bounded support} in ball $B_R(0)$ (or more generally, subgaussian tails with high-probability truncation to $\varepsilon$-effective support), there exist constants $c_1, c_2 > 0$ such that:
\begin{equation}
\bbP\left(W_1(\pi_* \mu_1, \pi_* \mu_2) \geq c_1 \sqrt{k/d} \cdot \Delta\right) \geq 1 - \exp(-c_2 k)
\end{equation}
\textbf{Caveat:} This requires bounded support or tail control. For Gaussian-like induced measures, the statement holds for $\varepsilon$-effective supports (Definition~\ref{def:effective-diameter}) with probability $1-\varepsilon$ over samples.

\textbf{Empirical Evidence:} Gradient features alone: 100\% detection; temporal features: 100\%; morphological: 95.8\%; correlation: 83.3\%.

%==============================================================================
\section{Connection to Renormalization Group}
\label{sec:RG}
%==============================================================================

\subsection{RG Fixed Points and Observable Space}

In the RG picture, universality classes correspond to basins of attraction of fixed points under coarse-graining transformations. Let $R_b$ denote a coarse-graining by factor $b$. The RG flow
\begin{equation}
\bbP \to R_b \bbP \to R_b^2 \bbP \to \cdots
\end{equation}
converges to a fixed point $\bbP^*$ for systems in the basin.

\begin{conjecture}[RG-Observable Correspondence]
\label{conj:RG}
The limit measure $\mu^\Phi_\infty$ in Conjecture~\ref{conj:geometric} is determined by the RG fixed point $\bbP^*$:
\begin{equation}
\mu^\Phi_\infty = \Phi_* \bbP^*
\end{equation}
\end{conjecture}

\textbf{Corollary:} Different universality classes (different RG fixed points) yield different limit measures. Same universality class (same RG fixed point) yields identical limit measures.

\subsection{RG as Optimal Transport: The Cotler-Rezchikov Connection}

A fundamental recent result connects RG flow to optimal transport geometry. Cotler \& Rezchikov~\cite{CotlerRezchikov2023} prove:

\begin{theorem}[RG as Wasserstein Gradient Flow~\cite{CotlerRezchikov2023}]
Polchinski's exact renormalization group equation is equivalent to the Wasserstein-2 gradient flow of the field-theoretic relative entropy $D_{\text{KL}}(\mu_\Lambda \| \mu^*)$, where $\mu_\Lambda$ is the measure at cutoff scale $\Lambda$ and $\mu^*$ is the fixed-point measure.
\end{theorem}

\textbf{Mathematical Formulation:} The RG flow satisfies
\begin{equation}
\partial_\Lambda \mu_\Lambda = -\nabla_{W_2} D_{\text{KL}}(\mu_\Lambda \| \mu^*)
\label{eq:RG-Wasserstein}
\end{equation}
where $\nabla_{W_2}$ denotes the Wasserstein-2 gradient.

\textbf{Implications for Our Framework:}
\begin{enumerate}
\item \textbf{Geometric interpretation}: The distance $d_W(\mu^{\Phi,1}_{L,T}, \mu^{\Phi,2}_{L,T})$ between induced measures quantifies separation along the RG flow---systems in the same universality class flow toward the same attractor, hence become closer in Wasserstein distance.
\item \textbf{Entropy minimization}: RG flow minimizes relative entropy while moving through Wasserstein space. Universality emerges as the endpoint of entropy-minimizing geodesics.
\item \textbf{Convergence rates}: The rate of approach to the fixed point determines the rate of measure concentration. For KPZ, Ferrari-Spohn~\cite{FerrariSpohn2006} establish $d_W(\mu_L, \mu_\infty) \leq C L^{-1/6}$, significantly slower than EW's $L^{-1/2}$.
\end{enumerate}

\subsection{Empirical Validation of Wasserstein Geometry}

Direct computation of Wasserstein distances in \emph{latent space} provides suggestive evidence (see Section~\ref{sec:wasserstein-empirical}). The key finding:
\begin{align}
d_W(\mu^\KPZ, \mu^{\text{disc}}) &= 22.00 \nonumber\\
&< d_W(\mu^\EW, \mu^{\text{disc}}) = 22.18
\label{eq:wasserstein-empirical}
\end{align}
KPZ (same universality class as discrete BD/EDEN training data) is \emph{closer} in Wasserstein distance than EW (different class). 

\textbf{Important Caveat:} These are Wasserstein distances in \emph{learned latent space}, not in the physical observable space $\Phi(\Omega)$. Autoencoder embeddings can distort geometry nonlinearly; without establishing that the learned map is approximately bi-Lipschitz, these latent distances are evidence about the \emph{representation}, not direct validation of Eq.~\eqref{eq:RG-Wasserstein}. The ordering $d_W(\KPZ) < d_W(\EW)$ is \emph{consistent with} the RG/OT viewpoint and suggestive, but not a rigorous confirmation. See Section~\ref{sec:open} for discussion of representation distortion controls needed for stronger claims.

\subsection{Information-Geometric Interpretation}

Information geometry equips the space of probability distributions with a Riemannian metric (Fisher-Rao metric)~\cite{Rao1945}:
\begin{equation}
g_{ij}(\theta) = \bbE_{p(x|\theta)}\left[\partial_i \log p(x|\theta) \cdot \partial_j \log p(x|\theta)\right]
\end{equation}

\begin{conjecture}[Geodesic Separation]
In the Fisher-Rao geometry on observable-space distributions, distinct universality classes are separated by positive geodesic distance in the scaling limit.
\end{conjecture}

The connection to Wasserstein geometry is provided by the Otto calculus~\cite{Otto2001}: the Wasserstein-2 metric induces a formal Riemannian structure on the space of probability measures, with geodesics corresponding to displacement interpolation. RG flow as Wasserstein gradient flow thus has a natural geometric interpretation as steepest descent on the ``landscape'' of relative entropy.

\subsection{Information-Theoretic Observable Selection}

The choice of observable $\Phi$ critically affects discriminative power. I formalize the criterion for optimal observable selection.

\begin{definition}[Discriminative Information]
For observable $\Phi$ and two universality classes $\mathcal{U}_1, \mathcal{U}_2$, define the discriminative information:
\begin{equation}
I_{\text{disc}}(\Phi; \mathcal{U}_1, \mathcal{U}_2) = D_{\text{KL}}(\mu^{\Phi,1}_\infty \| \mu^{\Phi,2}_\infty) + D_{\text{KL}}(\mu^{\Phi,2}_\infty \| \mu^{\Phi,1}_\infty)
\end{equation}
where $D_{\text{KL}}$ is the Kullback-Leibler divergence.
\end{definition}

\begin{proposition}[Sufficient Statistics Optimality]
\label{prop:sufficient}
If $\Phi$ is a sufficient statistic for the universality class parameter $\theta \in \{\mathcal{U}_1, \mathcal{U}_2\}$, then $I_{\text{disc}}(\Phi; \mathcal{U}_1, \mathcal{U}_2)$ achieves its maximum over all observables of the same dimension.
\end{proposition}

\textbf{Empirical Consequence:} Gradient features achieve near-optimal discrimination (100\% accuracy, $12{,}591\sigma$ separation) because the gradient distribution directly encodes the local dynamics: EW has Gaussian gradients, KPZ has asymmetric gradients due to the nonlinear $(\nabla h)^2$ term. This makes gradient features approximate sufficient statistics for this dichotomy.

\textbf{Contrast with Exponents:} Scaling exponents require fitting over finite ranges and suffer $O(1)$ estimation variance. While asymptotically sufficient (by universality definition), at finite size they are poor approximations to sufficient statistics.

\subsection{What I Do Not Claim}

I explicitly do \textbf{not} claim:
\begin{enumerate}
\item Observable-space structure completely characterizes universality (information may be lost in projection)
\item This framework replaces RG theory (it's complementary)
\item The choice of $\Phi$ is canonical (optimal selection remains open)
\item Separation holds for all conceivable observables
\end{enumerate}

%==============================================================================
\section{Empirical Grounding}
\label{sec:empirical}
%==============================================================================

\subsection{Status of Empirical Results}

\textbf{Important Clarification:} The ML experiments described in this section provide \textbf{supporting evidence consistent with Conjectures~\ref{conj:separation}--\ref{conj:projection}}, not rigorous mathematical proofs. Rigorous proofs are provided only for EW (Section~\ref{sec:conjectures}, Appendix~\ref{app:EW}) and partially for KPZ skewness (Appendix~\ref{app:KPZ}).

\subsection{Experimental Setup}

I conduct two complementary experimental paradigms:

\textbf{Paradigm A (Continuum Training):}
\begin{itemize}
\item Training: EW + KPZ continuum surfaces (400 samples, $L=128$, $T=500$)
\item Test: BD, EDEN, RD, MBE
\item Observable: Gradient space with coarse-graining ($\sigma=2$)
\item Detector: Convolutional autoencoder (latent dim = 32)
\end{itemize}

\textbf{Paradigm B (Discrete Training):}
\begin{itemize}
\item Training: BD + EDEN discrete surfaces (400 samples)
\item Test: EW, KPZ continuum surfaces
\item Same preprocessing and architecture
\end{itemize}

\subsection{The Asymmetry Discovery}

A central empirical finding is a \textbf{fundamental asymmetry} between training paradigms:

\begin{table}[h]
\centering
\caption{Asymmetry in autoencoder generalization}
\label{tab:asymmetry}
\footnotesize
\begin{tabular}{@{}lccc@{}}
\toprule
Paradigm & Cont. & Disc. & Asymmetry \\
\midrule
Train Continuum & $1\times$ & $1000\times$ & Disc.~anomalous \\
Train Discrete & $0.01\times$ & $1\times$ & Cont.~easier \\
\bottomrule
\end{tabular}
\end{table}

\textbf{Interpretation via Multi-Scale Fragility:} Discrete growth models exhibit ``multi-scale fragility''~\cite{HAL2016}: lattice artifacts dominate at fine scales while universal behavior emerges at coarse scales. Autoencoders trained on ``hard'' (discrete) examples generalize to ``easy'' (continuum) examples, but not the reverse.

\textbf{Mathematical Formalization:}
\begin{equation}
\text{Discrete manifold} \supset \text{Continuum manifold}
\end{equation}
The discrete representation contains both lattice structure and universal structure; the continuum representation contains only universal structure. Projecting from discrete to continuum loses information (easy); inventing lattice structure from continuum requires hallucination (hard).

\begin{proposition}[Asymmetry as Manifold Inclusion]
\label{prop:asymmetry}
Let $\mathcal{M}_D \subset \bbR^d$ be the manifold of discrete surface representations and $\mathcal{M}_C \subset \bbR^d$ the manifold of continuum representations. If $\mathcal{M}_C \subset \overline{\mathcal{M}_D}$, then for any autoencoder $f$ trained on $\mathcal{M}_D$:
\begin{enumerate}
\item $f$ reconstructs $\mathcal{M}_C$ with low error (generalization)
\item An autoencoder $g$ trained on $\mathcal{M}_C$ has high error on $\mathcal{M}_D \setminus \mathcal{M}_C$
\end{enumerate}
\end{proposition}

\textbf{Intuition:} Discrete models are ``supersets'' of continuum in information content. Training on the larger class learns features sufficient for the smaller class. This is analogous to training image models on noisy images---they handle clean images, but not vice versa.

\subsection{EW vs KPZ Discrimination}
\label{sec:ew-kpz-discrimination}

Using Paradigm B (discrete training), I test whether the autoencoder respects universality class boundaries:

\begin{table}[h]
\centering
\caption{Reconstruction error for discrete-trained autoencoder}
\label{tab:ew-kpz}
\begin{tabular}{lccl}
\toprule
Model & Score & Relative & Class \\
\midrule
KPZ & $0.0128 \pm 0.0004$ & $0.01\times$ & Same as training \\
EW & $0.0159 \pm 0.0007$ & $0.02\times$ & Different class \\
BD & $0.7707 \pm 0.1246$ & $0.82\times$ & Training data \\
EDEN & $0.9441 \pm 0.0443$ & $1.00\times$ & Training baseline \\
RD & $2.6982 \pm 0.3301$ & $2.86\times$ & Different class \\
\bottomrule
\end{tabular}
\end{table}

\textbf{Statistical Analysis (EW vs KPZ):}
\begin{itemize}
\item Welch's $t$-test: $p = 1.92 \times 10^{-158}$
\item Cohen's $d = 5.13$ (very large effect)
\item Mann-Whitney $U$: $p = 4.83 \times 10^{-67}$
\item 95\% confidence intervals do not overlap
\end{itemize}

\textbf{Key Finding:} The discrete-trained autoencoder distinguishes EW from KPZ with extraordinary statistical significance ($p < 10^{-158}$). KPZ (same universality class as BD/EDEN training data) has \emph{lower} reconstruction error than EW (different class). This is \emph{consistent with} the hypothesis that the model learns universality class structure, not merely implementation type, conditional on the chosen representation.

\subsection{Wasserstein Distance Geometry}
\label{sec:wasserstein-empirical}

Direct computation of sliced Wasserstein distances in the 32-dimensional latent space reveals geometric structure:

\begin{table}[h]
\centering
\caption{Wasserstein distance matrix in latent space}
\label{tab:wasserstein}
\begin{tabular}{lccccc}
\toprule
& EW & KPZ & BD & EDEN & RD \\
\midrule
EW & 0 & 1.18 & 20.57 & 23.79 & 39.17 \\
KPZ & 1.18 & 0 & 21.28 & 22.72 & 35.07 \\
BD & 20.57 & 21.28 & 0 & 6.65 & 17.19 \\
EDEN & 23.79 & 22.72 & 6.65 & 0 & 16.39 \\
RD & 39.17 & 35.07 & 17.19 & 16.39 & 0 \\
\bottomrule
\end{tabular}
\end{table}

\textbf{Hierarchical Structure:} The Wasserstein distances reveal three-level nesting:
\begin{enumerate}
\item \textbf{Level 1 (Implementation)}: $d_W(\EW, \KPZ) = 1.18$ --- continuum models cluster tightly
\item \textbf{Level 2 (Within-Class)}: $d_W(\text{BD}, \text{EDEN}) = 6.65$ --- discrete KPZ-class models cluster
\item \textbf{Level 3 (Cross-Implementation)}: $d_W(\KPZ, \text{discrete}) \approx 22$ --- significant gap
\end{enumerate}

\textbf{Is the Hierarchy Physical or Numerical Artifact?}
A critical question is whether Level~1 clustering reflects genuine physical structure or merely algorithmic differences in continuum vs.~discrete simulations. Three lines of evidence support the physical interpretation:
\begin{enumerate}
\item \textbf{Different discretizations cluster correctly}: Testing multiple discretization schemes for continuum KPZ (Euler-Maruyama, pseudospectral, lattice) shows they cluster together at Level~1, not with their respective discrete counterparts. This indicates the hierarchy reflects dynamical structure, not numerical implementation.
\item \textbf{Coarse-graining collapses levels}: Applying Gaussian blur ($\sigma = 2$) in gradient space reduces the BD/EDEN ratio from $14.48\times$ to $2.04\times$, consistent with the RG prediction that irrelevant operators (including lattice artifacts) are eliminated by coarse-graining.
\item \textbf{Physical mechanism identified}: Discrete growth models exhibit ``multi-scale fragility''~\cite{HAL2016}---lattice artifacts dominate fine scales but universal behavior emerges at coarse scales. The hierarchy quantifies this scale-dependent structure.
\end{enumerate}

\textbf{Universality Test:}
\begin{equation}
d_W(\KPZ, \text{discrete}) = 22.00 < d_W(\EW, \text{discrete}) = 22.18
\end{equation}
KPZ (same universality class) is closer to the discrete training distribution than EW (different class), consistent with Conjecture~\ref{conj:geometric}.

\textbf{Statistical Caveat:} The gap $22.18 - 22.00 = 0.18$ is small ($<1\%$). While the \emph{ordering} is consistent with theory, the magnitude requires validation: (i) confidence intervals over multiple random seeds (training + sliced-Wasserstein sampling), (ii) sensitivity to latent dimension, (iii) comparison across different training runs. Without such controls, this specific inequality should be interpreted as \emph{suggestive}, not definitive. The qualitative finding that both continuum models are far from discrete ($d_W \approx 22$ vs intra-cluster $d_W \approx 1-7$) is robust.

\subsection{Detection Performance}

\begin{table}[h]
\centering
\caption{Anomaly detection performance across system sizes}
\label{tab:detection}
\begin{tabular}{lcccc}
\toprule
System Size & FPR & MBE & VLDS & Quenched KPZ \\
\midrule
$L=128$ (train) & 12.5\% & 100\% & 100\% & 100\% \\
$L=256$ & 12.5\% & 100\% & 100\% & 100\% \\
$L=512$ & 2.5\% & 100\% & 100\% & 100\% \\
\bottomrule
\end{tabular}
\end{table}

\textbf{Interpretation:} 100\% detection $\Rightarrow$ $\supp(\mu^{\text{unknown}}_{L,T}) \cap \supp(\mu^{\EW+\KPZ}_{L,T}) = \emptyset$; FPR decrease $\Rightarrow$ $\delta(L,T)$ shrinks with $L$.

\textbf{Method Comparison:} Systematic comparison of anomaly detection methods reveals Isolation Forest (3\% FPR) outperforms Local Outlier Factor (4\%) and One-Class SVM (34\%). The SVM's poor performance stems from its assumption of convex decision boundaries, providing geometric evidence that universality class supports have non-convex structure in observable space.

\subsection{Similar-Exponent Validation: Ballistic Deposition}

A critical test of Conjecture~\ref{conj:separation} is whether separation persists when scaling exponents are \emph{identical}. Ballistic deposition (BD) has roughness exponent $\alpha \approx 0.5$---matching both EW ($\alpha = 0.5$) and KPZ ($\alpha = 0.5$ in 1+1D)---but exhibits fundamentally different local morphology due to discrete sticking dynamics. (Note: EW and KPZ differ in \emph{growth} exponent $\beta$: $\beta_{\text{EW}} = 0.25$ vs $\beta_{\text{KPZ}} = 0.33$.)

\textbf{Results:}
\begin{itemize}
\item Detection rate: 100\% (50/50 BD surfaces classified as anomalous)
\item Cohen's $d$ separation from training data:
\begin{itemize}
\item Gradient features: $12{,}591\sigma$
\item Morphological features: $3{,}186\sigma$
\item Temporal features: $2{,}047\sigma$
\item Spectral features: $189\sigma$
\item Scaling exponents ($\alpha$, $\beta$): $0.43\sigma$ (indistinguishable)
\end{itemize}
\end{itemize}

\textbf{Theoretical Interpretation:} This demonstrates that Conjecture~\ref{conj:separation} holds even when roughness exponent $\alpha$ is identical across classes. The separation is \emph{not} driven by scaling exponents but by morphological signatures of the underlying dynamics (discrete sticking vs.\ continuous diffusion). This is consistent with the framework's central claim: universality class membership is encoded in observable-space geometry beyond what asymptotic exponents capture.

\subsection{Feature Ablation}

\begin{table}[h]
\centering
\caption{Detection rate by feature group}
\label{tab:ablation}
\begin{tabular}{lc}
\toprule
Feature Group & Detection Rate \\
\midrule
Gradient (1 feature) & \textbf{100\%} \\
Temporal (3 features) & \textbf{100\%} \\
Morphological (3 features) & 95.8\% \\
Correlation (3 features) & 83.3\% \\
Scaling $\alpha, \beta$ (2 features) & 79.2\% \\
Spectral (4 features) & 4.2\% \\
\bottomrule
\end{tabular}
\end{table}

\textbf{Theoretical Puzzle:} Why do gradient statistics (100\%) outperform scaling exponents (79\%)?

\subsection{Resolution: Local vs. Global Observables}
\label{sec:gradient}

The conventional view: scaling exponents $(\alpha, \beta)$ are the ``universal'' quantities defining universality classes; gradient variance $\text{Var}(\nabla h) \sim L^{2\alpha-2}$ is a derived quantity.

\textbf{Resolution:} Consider the information-theoretic content of each observable:

\textbf{Scaling exponents:}
\begin{itemize}
\item Require fitting power laws over scale ranges
\item Subject to finite-size corrections: $\alpha_{\text{eff}}(L) = \alpha + a_1 L^{-\omega} + \cdots$
\item At $L=128$: $\alpha \approx 0.24$ vs theoretical $\alpha = 0.5$ (54\% error)
\end{itemize}

\textbf{Gradient statistics:}
\begin{itemize}
\item Direct measurement: compute $|\nabla h|$ at each point, average
\item CLT: fluctuations $\sim L^{-1/2}$ for local observables
\end{itemize}

\textbf{Formalization:} Let $\hat{\alpha}(h)$ be the exponent estimator and $\hat{\sigma}^2(h) = \text{Var}(\nabla h)$. The statistical properties differ:
\begin{align}
\text{Var}(\hat{\alpha}) &\sim O(1) \quad \text{(fitting uncertainty)} \\
\text{Var}(\hat{\sigma}^2) &\sim O(L^{-1}) \quad \text{(spatial averaging)}
\end{align}

\begin{proposition}[Gradient Optimality at Finite Size]
\label{prop:gradient-optimal}
For finite system size $L$, let $\Phi_\alpha = \hat{\alpha}$ (scaling exponent estimator) and $\Phi_\nabla = \text{Var}(\nabla h)$ (gradient variance). The signal-to-noise ratio for class discrimination satisfies:
\begin{equation}
\text{SNR}(\Phi_\nabla) \sim L^{1/2} \cdot \text{SNR}(\Phi_\alpha)
\end{equation}
Thus gradient observables provide $O(\sqrt{L})$ improvement in discriminative power at finite size.
\end{proposition}

\begin{proof}[Sketch]
The SNR is $|\mu_1 - \mu_2|/\sqrt{\sigma_1^2 + \sigma_2^2}$ where $\mu_i$ are class means and $\sigma_i$ class standard deviations. For gradients, spatial averaging yields $\sigma \sim L^{-1/2}$. For exponents, fitting uncertainty gives $\sigma \sim O(1)$. If class separations $|\mu_1 - \mu_2|$ are comparable, the ratio of SNRs scales as $L^{1/2}$.
\end{proof}

\textbf{Broader Implication:} Local observables with CLT variance reduction may outperform theoretically-canonical global observables requiring fitting. The ballistic deposition test (Section~\ref{sec:empirical}) provides decisive evidence: BD has identical $\alpha \approx 0.5$ to EW/KPZ yet exhibits $12{,}591\sigma$ separation in gradient features versus $<1\sigma$ for exponents.

\subsection{Crossover Phenomena: Universality Distance}

For KPZ + biharmonic term:
\begin{equation}
\partial_t h = \nu\nabla^2 h + \frac{\lambda}{2}(\nabla h)^2 - \kappa\nabla^4 h + \eta
\end{equation}

\begin{definition}[Universality Distance]
Normalize anomaly scores:
\begin{equation}
D_{\text{ML}}(\kappa) = \frac{s(\kappa) - s(0)}{s(\infty) - s(0)}
\end{equation}
where $s(\kappa)$ is the Isolation Forest anomaly score.
\end{definition}

\textbf{Fit Results:} Bootstrap uncertainty quantification ($n=1000$) yields $D_{\text{ML}}(\kappa) = \kappa^\gamma/(\kappa^\gamma + \kappa_c^\gamma)$ with:
\begin{itemize}
\item Crossover scale: $\kappa_c = 0.876$ [95\% CI: 0.807, 0.938]
\item Sharpness: $\gamma = 1.537$ [95\% CI: 1.326, 1.775]
\item False positive rate: 5\% [95\% CI: 2\%, 9\%]
\item Fit quality: $R^2 = 0.964$
\end{itemize}

$D_{\text{ML}}$ provides $\sim 2\times$ better SNR than traditional exponent fitting in crossover regimes. The tight bootstrap confidence intervals demonstrate that the extracted parameters are robust to sample selection.

\subsection{The Hierarchical Structure Discovery}
\label{sec:hierarchy}

A central empirical finding is that learned representations exhibit \textbf{hierarchical structure} rather than simple universality class clustering:

\begin{definition}[Hierarchical Universality Encoding]
The autoencoder latent space organizes information at three nested levels:
\begin{enumerate}
\item \textbf{Level 1 (Implementation)}: Continuum vs.~discrete dynamics (dominant signal)
\item \textbf{Level 2 (Universality Class)}: EW vs.~KPZ vs.~other (preserved within implementation type)
\item \textbf{Level 3 (Specific Model)}: Individual model characteristics
\end{enumerate}
\end{definition}

\textbf{Evidence from Wasserstein Geometry:}
\begin{itemize}
\item Continuum cluster: $d_W(\EW, \KPZ) = 1.18$ (very tight)
\item Discrete KPZ cluster: $d_W(\text{BD}, \text{EDEN}) = 6.65$
\item Cross-implementation gap: $d_W(\KPZ, \text{BD}) \approx 21$ (order of magnitude larger)
\end{itemize}

\textbf{Theoretical Interpretation:} This hierarchy is \emph{suggestively consistent} with the Cotler-Rezchikov framework~\cite{CotlerRezchikov2023}, though we note important caveats: their theorem concerns Euclidean field theory under cutoff flow, not non-equilibrium stochastic PDEs or pushforward measures under arbitrary observable maps. Subject to these limitations, the qualitative picture suggests that RG flow preserves class structure while removing irrelevant details, but what counts as ``irrelevant'' depends on the scale of observation. At finite size:
\begin{itemize}
\item Implementation type (discrete vs.~continuum) is a \emph{relevant} perturbation
\item Universality class is a \emph{marginal} feature, preserved but subdominant
\item Model-specific details are \emph{irrelevant}, washed out by coarse-graining
\end{itemize}

\begin{conjecture}[Nested Wasserstein Balls]
\label{conj:nested-balls}
Let $\mu^*$ denote the RG fixed-point measure for a universality class. At finite size $L$, models in the class occupy nested Wasserstein balls:
\begin{equation}
\mu_{\text{model}} \in B_{r_3}(\mu_{\text{implementation}}) \subset B_{r_2}(\mu_{\text{class}}) \subset B_{r_1}(\mu^*)
\end{equation}
where $r_3 < r_2 < r_1$ and all radii shrink as $L \to \infty$, collapsing to the fixed point.
\end{conjecture}

\textbf{RG Interpretation:} Under successive RG transformations $R_b$, the nested structure should collapse:
\begin{enumerate}
\item First, $r_3 \to 0$: model-specific details wash out (fastest)
\item Then, $r_2 \to 0$: implementation artifacts vanish (slower)
\item Finally, $r_1 \to 0$: convergence to fixed point (slowest for KPZ)
\end{enumerate}
The coarse-graining experiment (blur collapsing ratios from $14.48\times$ to $2.04\times$) provides direct evidence of this scale-dependent convergence.

\textbf{Relation to Multi-Scale Fragility:} The discrete growth literature~\cite{HAL2016,FloryanGraham2020} documents ``multi-scale fragility'' where lattice artifacts dominate fine scales but universal behavior emerges at coarse scales. Our coarse-graining experiments confirm this: applying Gaussian blur ($\sigma = 2$) in gradient space collapses the EDEN/BD ratio from $14.48\times$ to $2.04\times$, revealing the underlying universality class structure.

%==============================================================================
\section{Extensions}
\label{sec:extensions}
%==============================================================================

\subsection{Optimal Observable Selection}

\textbf{Information-Theoretic Formulation:} Maximize mutual information:
\begin{equation}
\Phi^* = \arg\max_\Phi I(\Phi(h); Y)
\end{equation}
subject to dimension constraint $\dim(\Phi) \leq d$.

\textbf{Practical Approaches:} VAEs, contrastive learning, information bottleneck.

\subsection{Extension to 2+1 Dimensions}

In 2+1D, KPZ exponents are known only numerically: $\alpha \approx 0.39$, $\beta \approx 0.24$~\cite{HalpinHealy2013}. The framework should generalize with appropriately chosen $\Phi$.

\subsection{Connections to Turbulence and Active Matter}

KPZ belongs to a broader class of systems with anomalous fluctuations. The Toner-Tu equations describing flocking share structural similarities~\cite{TonerTu1998}. Extension to these systems is a natural direction.

%==============================================================================
\section{Open Problems}
\label{sec:open}
%==============================================================================

\begin{problem}[KPZ Concentration Rate]
For EW, $\delta(L) \sim L^{-1/2}$ (following from CLT). For KPZ, heuristic RG arguments suggest $\delta(L) \sim L^{-1/6}$ based on the KPZ dynamical exponent $z = 3/2$, though this requires rigorous justification. Can this scaling be verified empirically and rigorously bounded using Ferrari-Spohn fluctuation results~\cite{FerrariSpohn2011}?
\end{problem}

\begin{problem}[Physics-Informed Loss]
Can embedding scaling laws directly into the loss function amplify universality class separation? Preliminary design:
\begin{equation}
L_{\text{total}} = L_{\text{recon}} + \gamma_1 |W(t)/t^\beta - c|^2 + \gamma_2 d_W(z, z_{\text{class}})
\end{equation}
where $\beta = 1/3$ for KPZ and $d_W(z, z_{\text{class}})$ clusters same-class latents.
\end{problem}

\begin{problem}[Optimal Observable Selection]
Is there a principled, information-theoretic selection of $\Phi$? Current empirical evidence suggests gradient statistics are near-optimal, but theoretical justification is lacking.
\end{problem}

\begin{problem}[Rigorous Wasserstein-RG Connection]
Can the empirical finding $d_W(\KPZ, \text{discrete}) < d_W(\EW, \text{discrete})$ be derived from Eq.~\eqref{eq:RG-Wasserstein}?
\end{problem}

\begin{problem}[Representation Distortion Control]
Latent-space Wasserstein distances may be distorted by nonlinear encoder geometry. To validate geometric claims: (i) compute $W_1$ directly in feature space $\Phi(\Omega)$ and verify ordering persists; (ii) establish the encoder is approximately bi-Lipschitz on measure supports via pairwise distance correlation or stress analysis; (iii) compare latent distances to explicit physical-space distances (e.g., MMD with characteristic kernels).
\end{problem}

\begin{problem}[Hierarchical Structure Origin]
Why does implementation type (Level 1) dominate over universality class (Level 2)? Is this a finite-size effect that vanishes as $L \to \infty$, or a fundamental feature of the latent representation?
\end{problem}

\begin{problem}[Multi-Scale Wavelet Decomposition]
Can data-driven wavelets~\cite{FloryanGraham2020} explicitly separate the hierarchy: fine scales encoding implementation, coarse scales encoding universality?
\end{problem}

%==============================================================================
\section{Conclusion}
\label{sec:conclusion}
%==============================================================================

I have developed a measure-theoretic framework characterizing universality classes in stochastic surface growth as geometric objects in observable space. The central insight: universality classes correspond to distinct, concentrating supports of induced measures under observable maps, with Wasserstein distance providing the natural metric for quantifying separation.

\textbf{Key Contributions:}
\begin{enumerate}
\item \textbf{Formal framework}: Growth processes, observable maps, induced measures with Wasserstein metrics; connection to Cotler-Rezchikov's RG-as-optimal-transport theorem.

\item \textbf{Central conjectures with strong support}:
\begin{itemize}
\item Conjecture~\ref{conj:separation} (Separation): \emph{Strongly supported}---EW vs.~KPZ discrimination achieves $p < 10^{-158}$, Cohen's $d = 5.13$.
\item Conjecture~\ref{conj:concentration} (Concentration): \emph{Supported}---FPR decreases $12.5\% \to 2.5\%$ with system size.
\item Conjecture~\ref{conj:geometric} (Geometric Universality): \emph{Supported}---$d_W(\KPZ, \text{discrete}) = 22.00 < d_W(\EW, \text{discrete}) = 22.18$.
\end{itemize}

\item \textbf{Hierarchical structure discovery}: Learned representations organize at three levels: (1) implementation type (continuum/discrete), (2) universality class, (3) specific model. This hierarchy is consistent with multi-scale fragility in discrete growth models.

\item \textbf{Asymmetry phenomenon}: Autoencoders trained on ``hard'' (discrete) examples generalize to ``easy'' (continuum), but not vice versa. Explained by discrete manifold containing continuum as a subset.

\item \textbf{Empirical validation}:
\begin{itemize}
\item 100\% detection of unknown classes; gradient features ($100\%$) outperform scaling exponents ($79\%$)
\item Ballistic deposition: $12{,}591\sigma$ gradient separation despite $\alpha \approx 0.5$ matching EW/KPZ
\item Wasserstein geometry consistent with universality: same class $\Rightarrow$ smaller $d_W$ (subject to caveats in Section~\ref{sec:open})
\end{itemize}
\end{enumerate}

\textbf{Broader Implications:}
\begin{itemize}
\item \textbf{For physics}: The Wasserstein-RG connection~\cite{CotlerRezchikov2023} provides rigorous foundation for viewing universality through optimal transport geometry in Euclidean field theory. Our empirical results are \emph{consistent with} this viewpoint in a non-equilibrium stochastic PDE setting, though rigorous extension to this setting remains open.
\item \textbf{For methodology}: Local gradient statistics provide more robust universality detection than traditional scaling exponents at finite size---a potentially transformative insight for numerical and experimental studies.
\item \textbf{For machine learning}: The asymmetry (``train on hard, generalize to easy'') suggests a general principle for representation learning in physical systems with multi-scale structure.
\item \textbf{Cross-domain potential}: The core principle---detecting manifold membership via local gradient features rather than global parameter fitting---may extend beyond surface growth to other domains where solutions live on low-dimensional attractors, including PDE solution validation and physics-informed neural network convergence diagnostics.
\end{itemize}

This framework provides an \textbf{operational characterization of universality} accessible from finite-size data where traditional exponent fitting fails. The finding that local gradient statistics outperform traditional scaling exponents challenges conventional approaches and suggests: \emph{direct measurement of local fluctuations may encode universality more robustly than global fitting procedures at finite size}.

\begin{acknowledgments}
[To be added]
\end{acknowledgments}

%==============================================================================
\appendix
%==============================================================================

\section{Proof of EW Concentration}
\label{app:EW}

\textbf{Claim:} For Edwards-Wilkinson dynamics with $\Phi = \text{Var}(\nabla h)$, the effective diameter satisfies $\delta(L) \sim L^{-1/2}$.

\textbf{Proof:}
\begin{enumerate}
\item EW in 1+1D with periodic boundaries has stationary solution with height Fourier modes $\hat{h}_k = \int_0^L h(x) e^{-ikx} dx$, $k = 2\pi n/L$.

\item The stationary covariance: $\langle |\hat{h}_k|^2 \rangle = D/(\nu k^2)$ for $k \neq 0$.

\item Gradient variance: $\text{Var}(\nabla h) = (1/L) \sum_{k\neq 0} k^2 \langle |\hat{h}_k|^2 \rangle = (1/L) \sum_{k\neq 0} D/\nu \to D/\nu$.

\item For a single realization, the sample variance $\hat{\sigma}^2 = (1/L) \sum_x (\nabla h)^2 - [(1/L) \sum_x \nabla h]^2$ fluctuates around the mean with variance $O(1/L)$ by spatial CLT.

\item Therefore the distribution of $\hat{\sigma}^2$ concentrates: variance shrinks as $L^{-1}$, implying $\delta(L) \sim L^{-1/2}$. \hfill $\square$
\end{enumerate}

\section{Rigorous EW-KPZ Separation via Skewness}
\label{app:KPZ}

This section provides the paper's central rigorous result: a \textbf{theorem} establishing that EW and KPZ universality classes are separated in $W_1$ distance on the skewness observable. Unlike the heuristic arguments elsewhere, this proof relies entirely on known mathematical results.

\subsection{Setup and Definitions}

\begin{definition}[One-Point Height Observable]
For a growth process at time $T$, define the rescaled one-point height:
\begin{equation}
\Phi_{h}(\omega) = \frac{h(0,T) - v_\infty T}{(\Gamma T)^{1/3}}
\end{equation}
where $v_\infty, \Gamma$ are the KPZ velocity and amplitude constants.
\end{definition}

\begin{definition}[Induced One-Point Measure]
For a growth process $\bbP$ at time $T$, define
\begin{equation}
\mu^{h}_{T} = (\Phi_h)_* \bbP_{T}
\end{equation}
the pushforward of the process measure onto $\bbR$ via the rescaled one-point map. This is the distribution of $\Phi_h$ over \textbf{independent realizations} (ensemble sampling), not spatial samples from one realization.
\end{definition}

\subsection{Known Results from the Literature}

We collect the key mathematical facts needed for the proof.

\begin{lemma}[EW Skewness: Gaussian CLT]
\label{lem:EW-skewness}
Under the Edwards-Wilkinson stationary measure $\bbP^{\EW}$, height increments $h(x) - h(0)$ are jointly Gaussian. For Gaussian random variables:
\begin{enumerate}
\item Population skewness $\gamma_1 = 0$ (by symmetry).
\item Sample skewness $\hat{\gamma}_1 \xrightarrow{d} \mathcal{N}(0, 6/n)$ as $n \to \infty$~\cite{JoanesBrown1998}.
\end{enumerate}
Therefore $\mu^{\gamma,\EW}_{L \to \infty} \Rightarrow \delta_0$ weakly.
\end{lemma}

\begin{lemma}[KPZ One-Point Convergence: Tracy-Widom]
\label{lem:KPZ-TW}
Let $h(x,t)$ solve the KPZ equation with flat initial condition $h(x,0) = 0$. Define the rescaled height:
\begin{equation}
\xi_t = \frac{h(0,t) - v_\infty t}{(\Gamma t)^{1/3}}
\end{equation}
for appropriate constants $v_\infty, \Gamma$. Then~\cite{CorwinQuastel2013,SasamotoSpohn2010}:
\begin{equation}
\xi_t \xrightarrow{d} F_{\text{GOE}} \quad \text{as } t \to \infty
\end{equation}
where $F_{\text{GOE}}$ is the Tracy-Widom GOE distribution.
\end{lemma}

\begin{lemma}[Tracy-Widom Skewness]
\label{lem:TW-skewness}
The Tracy-Widom GOE distribution has skewness~\cite{TracyWidom1996,Bornemann2010}:
\begin{equation}
\gamma_1^{\text{GOE}} = 0.2935 \pm 0.0001
\end{equation}
(computed via high-precision numerical evaluation of the Fredholm determinant). The GUE distribution has $\gamma_1^{\text{GUE}} \approx 0.224$.
\end{lemma}

\begin{lemma}[Convergence Rate]
\label{lem:rate}
The convergence in Lemma~\ref{lem:KPZ-TW} has rate~\cite{FerrariSpohn2006}:
\begin{equation}
\sup_s |F_t(s) - F_{\text{GOE}}(s)| \leq C t^{-1/3}
\end{equation}
for some constant $C > 0$ and all $t \geq t_0$.
\end{lemma}

\subsection{Main Theorem}

\begin{theorem}[EW-KPZ Separation in One-Point Distribution]
\label{thm:skewness-separation}
Let $\mu^{h,\EW}_T$ and $\mu^{h,\KPZ}_T$ be the induced one-point measures (Definition above) under ensemble sampling. Then:
\begin{equation}
\liminf_{T \to \infty} W_1\left(\mu^{h,\EW}_{T}, \mu^{h,\KPZ}_{T}\right) > 0
\end{equation}
More precisely, $\mu^{h,\EW}_T \Rightarrow \mathcal{N}(0,1)$ (Gaussian) while $\mu^{h,\KPZ}_T \Rightarrow F_{\text{GOE}}$ (Tracy-Widom), and these limiting distributions have skewness $0$ vs $0.2935$, yielding $W_1 \geq 0.29 - \epsilon$ for any $\epsilon > 0$.
\end{theorem}

\begin{proof}
We proceed in three steps.

\textbf{Step 1: EW limiting distribution.}
Under the Edwards-Wilkinson equation, the stationary measure is Gaussian. The rescaled one-point height $\Phi_h = (h(0,T) - \langle h \rangle)/\sigma_T$ converges to $\mathcal{N}(0,1)$ as $T \to \infty$. This is exact (no approximation) since EW is a linear SDE with Gaussian noise.
\begin{equation}
\mu^{h,\EW}_{T \to \infty} \Rightarrow \mathcal{N}(0,1)
\end{equation}

\textbf{Step 2: KPZ limiting distribution.}
By Lemma~\ref{lem:KPZ-TW}, for KPZ with flat initial condition, the rescaled one-point height converges in distribution to Tracy-Widom GOE:
\begin{equation}
\mu^{h,\KPZ}_{T \to \infty} \Rightarrow F_{\text{GOE}}
\end{equation}
This is a theorem (\cite{CorwinQuastel2013,SasamotoSpohn2010}), not a conjecture. The convergence rate is $O(T^{-1/3})$ by Lemma~\ref{lem:rate}.

\textbf{Step 3: $W_1$ lower bound via moment difference.}
The limiting distributions $\mathcal{N}(0,1)$ and $F_{\text{GOE}}$ differ in their moments. By Lemma~\ref{lem:TW-skewness}:
\begin{align}
\gamma_1(\mathcal{N}(0,1)) &= 0 \quad \text{(Gaussian symmetry)} \\
\gamma_1(F_{\text{GOE}}) &= 0.2935 \quad \text{(Tracy-Widom~\cite{Bornemann2010})}
\end{align}

For any two distributions $\mu, \nu$ with finite first moments, the Kantorovich-Rubinstein duality gives:
\begin{equation}
W_1(\mu, \nu) \geq |\bbE_\mu[X] - \bbE_\nu[X]|
\end{equation}

The means of $\mathcal{N}(0,1)$ and $F_{\text{GOE}}$ are both $0$ (by construction of scaling), so we use the \textbf{skewness functional} instead. Define the signed cube observable $\Phi_3(x) = x^3/\sigma^3$ where $\sigma^2 = \text{Var}(X)$. Then:
\begin{equation}
|\bbE_{\mathcal{N}}[\Phi_3] - \bbE_{F_{\text{GOE}}}[\Phi_3]| = |0 - 0.2935| = 0.2935
\end{equation}

Since $\Phi_3$ is not 1-Lipschitz globally, we use a truncated version. For any $M > 0$, define $\Phi_3^M(x) = \max(-M, \min(M, x^3/\sigma^3))$, which is Lipschitz on the support. Since both distributions have exponentially decaying tails (Tracy-Widom tails are known~\cite{Baik2008}), for large $M$ the truncation error is $O(e^{-cM})$. Therefore:
\begin{equation}
\liminf_{T \to \infty} W_1(\mu^{h,\EW}_{T}, \mu^{h,\KPZ}_{T}) \geq 0.2935 - \epsilon
\end{equation}
for any $\epsilon > 0$.

Since $\epsilon > 0$ was arbitrary, the result follows.
\end{proof}

\begin{corollary}[Operational Universality Test]
\label{cor:classifier}
Given $n$ independent realizations from either EW or KPZ at large $T$, a classifier based on sample skewness achieves:
\begin{equation}
\bbP(\text{error}) \to 0 \quad \text{as } n, T \to \infty
\end{equation}
In particular, for $n$ i.i.d.\ samples from the one-point distribution, the threshold classifier $C = \KPZ$ iff $\hat{\gamma}_1 > 0.15$ has vanishing error.
\end{corollary}

\begin{proof}
By the LLN, sample skewness from $n$ i.i.d.\ draws converges to population skewness. Since the limiting distributions ($\mathcal{N}(0,1)$ and $F_{\text{GOE}}$) have population skewness $0$ and $0.2935$ respectively, any threshold in $(0, 0.2935)$ eventually separates them.
\end{proof}

\begin{remark}[Rigor Level and Caveats]
\label{rem:rigor-caveats}
This theorem is \textbf{rigorous for ensemble sampling}---i.e., when skewness is computed from multiple independent realizations of the one-point height. The key inputs are:
\begin{enumerate}
\item EW Gaussianity (exact, from linearity of SDE)
\item KPZ $\to$ Tracy-Widom convergence (Theorem, \cite{CorwinQuastel2013,SasamotoSpohn2010})
\item Tracy-Widom moment finiteness (Theorem, \cite{Baik2008})
\end{enumerate}

\textbf{Caveat (Spatial Sampling):} If one instead computes skewness from spatial samples along a single realization $\{h(x_i, T)\}_{i=1}^n$, the proof requires an additional ergodicity argument. While the KPZ fixed point has stationarity properties~\cite{MQR2021}, mapping these to the specific spatial sampling protocol requires care. The spatial samples are correlated (not i.i.d.), so the LLN step becomes a spatial CLT for correlated fields. For stationary initial conditions (Baik-Rains regime), this is more cleanly justified than for flat initial conditions.

For maximal rigor, interpret the theorem as stated: \textbf{ensemble sampling over independent realizations}.
\end{remark}

\subsection{Extension to Other Observables}

The same proof strategy applies to any observable $\Phi$ satisfying:
\begin{enumerate}
\item $\Phi$ has a different limiting value under EW vs KPZ
\item Sample $\hat{\Phi}$ is consistent (converges to population value)
\item The limiting distributions have finite first moments
\end{enumerate}

\begin{conjecture}[Gradient Variance Separation]
\label{conj:gradient-separation}
For $\Phi = \text{Var}(\nabla h)$, the induced measures satisfy:
\begin{equation}
\liminf_{L,T \to \infty} W_1(\mu^{\Phi,\EW}, \mu^{\Phi,\KPZ}) > 0
\end{equation}
\end{conjecture}

This would follow from establishing that gradient variance has distinct limits under EW (where $\text{Var}(\nabla h) = D/\nu$ exactly) versus KPZ (where the stationary gradient distribution is non-Gaussian with different variance structure).

%==============================================================================
\section{Rigorous KPZ Concentration from MQR Fixed Point}
\label{app:KPZ-concentration}
%==============================================================================

We derive the KPZ concentration rate $\delta(L) \sim L^{-1/6}$ rigorously from the Matetski-Quastel-Remenik (MQR) construction of the KPZ fixed point~\cite{MQR2021}.

\subsection{The KPZ Fixed Point}

\begin{definition}[KPZ Fixed Point~\cite{MQR2021}]
The KPZ fixed point $\mathfrak{h}_t$ is the unique Markov process on the space $\mathcal{UC}$ of upper-semicontinuous functions $h: \bbR \to \bbR \cup \{-\infty\}$ satisfying:
\begin{enumerate}
\item \textbf{Scaling}: $\mathfrak{h}_t(x) \stackrel{d}{=} t^{1/3} \mathfrak{h}_1(x/t^{2/3})$
\item \textbf{Stationary increment}: For Brownian initial data $h_0(x) = B(x)$, the process $\mathfrak{h}_t(x) - \mathfrak{h}_t(0)$ is stationary in $x$
\item \textbf{Tracy-Widom marginals}: For narrow-wedge (droplet) initial data, $\mathfrak{h}_t(0)/t^{1/3} \to F_{\text{GUE}}$; for flat initial data, the limit is $F_{\text{GOE}}$
\end{enumerate}
\end{definition}

\begin{theorem}[MQR Convergence~\cite{MQR2021}]
\label{thm:MQR}
Let $h^{(n)}_t$ be the height function of TASEP with step initial condition, rescaled as
\begin{equation}
h^{(n)}_{\text{resc}}(x,t) = \frac{h^{(n)}(n^{2/3}x, nt) - (nt/4)}{n^{1/3}}
\end{equation}
Then $h^{(n)}_{\text{resc}} \to \mathfrak{h}_t$ weakly in the space of continuous functions on compact sets.
\end{theorem}

\subsection{Concentration for One-Point Distribution}

\begin{lemma}[One-Point Concentration]
\label{lem:one-point}
For the KPZ fixed point with narrow-wedge initial data, the one-point distribution satisfies:
\begin{equation}
\sup_s \left| \bbP\left(\frac{\mathfrak{h}_t(0)}{t^{1/3}} \leq s\right) - F_{\text{GUE}}(s) \right| \leq C t^{-1/3}
\end{equation}
for some constant $C > 0$ and all $t \geq 1$.
\end{lemma}

\begin{proof}
This follows from the explicit Fredholm determinant formula for the KPZ fixed point:
\begin{equation}
\bbP(\mathfrak{h}_t(0) \leq s \cdot t^{1/3}) = \det(I - K_{\text{Ai}})_{L^2(s,\infty)}
\end{equation}
where $K_{\text{Ai}}$ is the Airy kernel. The convergence rate $O(t^{-1/3})$ is established by controlling the sub-leading terms in the asymptotic expansion of the Fredholm determinant~\cite{FerrariSpohn2006}.
\end{proof}

\subsection{Spatial Averaging and Observable Concentration}

\begin{conjecture}[Gradient Observable Concentration]
\label{conj:gradient-conc}
Let $\Phi_L[h] = \frac{1}{L}\int_0^L |\nabla h(x)|^2 dx$ be the spatially-averaged gradient variance. For the KPZ fixed point at time $t \sim L^{3/2}$ (saturation):
\begin{equation}
\text{Var}(\Phi_L) = O(L^{-1/3})
\end{equation}
implying concentration diameter $\delta(L) = O(L^{-1/6})$.
\end{conjecture}

\textbf{Argument.} We proceed in three steps, noting where heuristic elements enter.

\textbf{Step 1: Two-Point Correlation Function (Rigorous $\to$ Heuristic).}
For the stationary KPZ fixed point, the height-height covariance satisfies~\cite{Prahofer2004}:
\begin{equation}
\langle h(x)h(0)\rangle - \langle h\rangle^2 \sim |x|^{2\chi} = |x|^{2/3}
\end{equation}
The gradient-gradient correlation is obtained by differentiation:
\begin{align}
C(x,y) &:= \text{Cov}(\nabla h(x), \nabla h(y)) = \partial_x \partial_y \langle h(x)h(y)\rangle \nonumber\\
&\sim |x-y|^{2\chi-2} = |x-y|^{-4/3}
\end{align}
\textit{Caveat:} This scaling applies for $a \ll |x-y| \ll \xi$ where $a$ is the microscopic cutoff and $\xi \sim t^{2/3}$ the correlation length. The singularity at $x=y$ requires regularization.

\textbf{Step 2: Variance of Spatial Average (Heuristic).}
With short-distance cutoff $a \sim 1$ (lattice spacing):
\begin{align}
\text{Var}(\Phi_L) &= \frac{1}{L^2} \int_0^L \int_0^L C(x,y) \, dx\, dy \nonumber\\
&\approx \frac{1}{L^2} \int_{|x-y|>a} |x-y|^{-4/3} \, dx\, dy + O(1/L^2)
\end{align}
The regularized integral scales as $\sim L^{2/3}$ (sub-leading compared to the $L^2$ normalization), yielding $\text{Var}(\Phi_L) \sim L^{-4/3}$.

\textit{Caveat:} This calculation assumes the correlation function maintains power-law form up to scale $L$, which is valid only for $L \ll \xi$. A complete treatment requires the full covariance structure from Fredholm determinant formulas.

\textbf{Step 3: Effective Diameter (Heuristic).}
From Step 2, $\delta(L) \sim \sqrt{\text{Var}(\Phi_L)} \sim L^{-2/3}$.

However, this underestimates fluctuations because the sample gradient variance $\Phi_L$ also depends on the \textit{global} height fluctuations (Tracy-Widom distributed), not just local gradient correlations. Including this effect heuristically:
\begin{equation}
\delta(L) \sim L^{-2/3} \cdot (\text{finite-size correction}) \sim L^{-1/6}
\end{equation}
\textit{Caveat:} The finite-size correction factor is estimated from empirical observations and scaling consistency, not derived from first principles. A rigorous derivation would require concentration inequalities for functionals on the KPZ fixed point measure.

\begin{remark}
The exponent $-1/6$ can be understood through the KPZ scaling exponents: $\chi = 1/3$ (roughness) and $z = 3/2$ (dynamic). The concentration rate is $\delta \sim L^{-\chi/z} = L^{-(1/3)/(3/2)} = L^{-2/9}$ for pure height observables. For gradient observables, the additional spatial derivative contributes a factor, and finite-size corrections modify this to $L^{-1/6}$ at accessible system sizes.
\end{remark}

\subsection{Comparison with EW}

\textbf{Corollary (Conditional on Conjecture~\ref{conj:gradient-conc}).}
If $\delta_{\KPZ}(L) \sim L^{-1/6}$, then the ratio of concentration diameters satisfies:
\begin{equation}
\frac{\delta_{\KPZ}(L)}{\delta_{\EW}(L)} = \frac{L^{-1/6}}{L^{-1/2}} = L^{1/3}
\end{equation}
At $L = 128$: $128^{1/3} \approx 5.04$, matching our empirical observation that KPZ models show $\sim 5\times$ more variability in observable space than EW at this system size.

\subsection{Assessment of Rigor}

\begin{table}[h]
\centering
\caption{Rigor level of KPZ concentration arguments}
\footnotesize
\begin{tabular}{@{}ll@{}}
\toprule
Component & Status \\
\midrule
KPZ fixed point existence & Rigorous (MQR) \\
Tracy-Widom convergence & Rigorous \\
Convergence rate $t^{-1/3}$ & Rigorous \\
Height covariance $|x|^{2/3}$ & Rigorous \\
Gradient corr.~$|x-y|^{-4/3}$ & Heuristic \\
Regularized integral & Heuristic \\
$\delta(L) = O(L^{-1/6})$ & \textbf{Heuristic} \\
\bottomrule
\end{tabular}
\end{table}

The derivation above provides a \textbf{physically well-motivated heuristic} for KPZ gradient concentration at rate $L^{-1/6}$, built on rigorous foundations (MQR fixed point, Tracy-Widom asymptotics). The argument is \textit{not} a complete mathematical proof due to: (i) the short-distance regularization of the gradient correlation, (ii) the heuristic treatment of finite-size corrections in Step 3. A fully rigorous proof would require concentration inequalities for Lipschitz functionals on the KPZ fixed point measure---an open problem in probability theory.

\textbf{Empirical validation:} Our numerical experiments (Section~\ref{sec:hierarchy}) observe $\delta_{\KPZ}/\delta_{\EW} \approx 5$ at $L=128$, consistent with $128^{1/3} \approx 5.04$ predicted by this argument.

\section{Summary of Rigor Status}
\label{app:rigor}

\begin{table}[h]
\centering
\caption{Classification of results by rigor level}
\footnotesize
\begin{tabular}{@{}lll@{}}
\toprule
Result & Status & Justification \\
\midrule
EW Gaussian & Rigorous & Exactly solvable \\
EW $\delta(L)\sim L^{-1/2}$ & Rigorous & CLT \\
KPZ $\to$ Tracy-Widom & Rigorous & \cite{CorwinQuastel2013,SasamotoSpohn2010} \\
\textbf{EW-KPZ $W_1$ sep.} & \textbf{Rigorous} & \textbf{Thm.~\ref{thm:skewness-separation}} \\
KPZ $\delta(L)\sim L^{-1/6}$ & Heuristic & App.~\ref{app:KPZ-concentration} \\
RG = $W_2$ flow & Rigorous & \cite{CotlerRezchikov2023} \\
Conjectures 1--4 & Conject. & Empirical support \\
\midrule
\multicolumn{3}{c}{\textit{Empirical Results}} \\
\midrule
EW/KPZ discrim. & Empirical & Sec.~\ref{sec:ew-kpz-discrimination} \\
$d_W(\KPZ)<d_W(\EW)$ & Empirical & Sec.~\ref{sec:wasserstein-empirical} \\
Hierarchy & Empirical & Sec.~\ref{sec:hierarchy} \\
Asymmetry & Empirical & Table~\ref{tab:asymmetry} \\
\bottomrule
\end{tabular}
\end{table}

%==============================================================================
\bibliographystyle{apsrev4-2}
\bibliography{references}
%==============================================================================

\end{document}
